\begin{apartats}

\apartat[1,25]{0,75}
Comprova que $F(x)=\ln(x^2+4)$ és una primitiva de $f(x)=\dfrac{2x}{x^2+4}$. Quina és la primitiva
de $f$ que val 0 a $x=0$?

\begin{solucio}
$F'(x)=\dfrac{1}{x^2+4}\cdot2x=\dfrac{2x}{x^2+4}=f(x)$ per a tot $x$.\\
Les primitives de $f$ són $\ln(x^2+4)+k$. A $x=0$ valen $\ln4+k$, que és 0 si $k=-\ln4$: la primitiva
buscada és $\ln(x^2+4)-\ln4=\ln\dfrac{x^2+4}{4}$.
\end{solucio}

\begin{tria}{primitiva-condicio}
\itemtria{exponencial}{1}{1,25}
Troba la primitiva de $f(x)=e^{2x}+\dfrac1x$, per a $x>0$, que compleix $F(1)=\dfrac{e^2}{2}$.

\begin{solucio}
$F(x)=\dfrac{e^{2x}}{2}+\ln x+k$. Com que $F(1)=\dfrac{e^2}{2}+0+k$, la condició dona $k=0$:
$F(x)=\dfrac{e^{2x}}{2}+\ln x$.
\end{solucio}

\itemtria{velocitat}{1}{1,25}
Un mòbil es desplaça en línia recta amb velocitat $v(t)=3t^2-2t$ (en m/s). Sabent que la posició és una
primitiva de la velocitat, i que a l'instant $t=0$ el mòbil és a $s(0)=5$\,m, troba la posició $s(t)$ i
on és a l'instant $t=2$.

\begin{solucio}
$s(t)=t^3-t^2+k$, i $s(0)=k=5$: $s(t)=t^3-t^2+5$. A l'instant $t=2$, $s(2)=8-4+5=9$\,m.
\end{solucio}
\end{tria}

\begin{nomesllarg}
\apartat{0,75}
Si $f$ és un polinomi de grau 3, de quin grau són les seves primitives? Troba la primitiva de
$f(x)=4x^3-6x$ que s'anul·la a $x=1$.

\begin{solucio}
De grau 4: integrar $x^n$ dona $\frac{x^{n+1}}{n+1}$, que té un grau més (i derivar en treu un).\\
$F(x)=x^4-3x^2+k$, i $F(1)=1-3+k=0$ dona $k=2$: $F(x)=x^4-3x^2+2$.
\end{solucio}
\end{nomesllarg}

\end{apartats}
